\section{Citadel 2.0: Improving Performance}

Our first baseline successfully achieves security.
In this section, we show several mechanisms to improve its performance.
First. fixed LLC partitioning \'{a} la MI6 destroys performance for our 2-core machine.
We show a reconfigurable LLC partitioning scheme and a fine grain flushing mechanism.
Furthermore, disabling speculative shared memory accesses makes software requirements especially light, but provides limited performances.
We show a relaxed speculative model where we enable pipelined accesses to shared memory but disable predictors and hard-code them to $pc+4$.
When this mode is enabled, the program needs to verify more requirements, i.e., that no secret is leaked when speculating $pc+4$.
We show that this new mode improves performance for code segments with burst shared memory accesses, but it hampers performances in the general case.
We observe that this is due to the TLB being trashed by misspeculations across functions context. 
This is an effect that was studied (in a nonsecurity context) in the early 2000s\cite{onur1,onur2,onur3}: misspeculations can induce large overhead not only due to the delay in executing the right path, large delays are induced due to the trashing of the TLBs.
Finally, we show how insights from building the SM make it possible to securely optimize context switch and enable frequent interrupts.

\subsection{Evaluation-Guided Optimisations}

Our secure baseline runs on FPGA with a Linux environment (see section \ref{sec:eval} for our full configuration) which makes it possible to analyses performance bottlenecks on real programs.
For this section, we use an elliptic curves (Ed25519) cryptographic library as a running example (see section \ref{sec:evalcrypto} for more details).
At a high level, our enclave interactively receives messages to sign from a benchtest program using shared memory.
Out-of-the box, we see an underwhelming overhead of 120\% overhead compared to directly (and insecurely) linking the cryptographic library to our benchtest program.

\subsection{Reconfigurable LLC Partitioning}
\label{sec:dynamiccache}
An obvious bottleneck in our design is cache partitioning.
Sanctum and MI6's mechanism (i.e. dividing the LLC in 64 equal size slices to match the 64 DRAM regions) is a terrible idea for our system.
Our 2-core prototype has a 1MB LLC (see \ref{tab:mi6param}), which divided by 64 gives us 16KB LLC slices.
This is smaller than the 64KB L1 D+I caches.
Because our caches are inclusive, this means an enclave or a program that resides in one memory region cannot even fully use the L1 caches.
MI6 circumvented that issue by assuming a 16-core design with a large 16MB LLC.

We build a mechanism for the OS to reconfigure the LLC partitions to improve enclave and other program's performance.
This costly operation is only required when changing the LLC layout, which only happens before launching an application or an enclave with an important memory footprint.
The OS requests a reconfiguration through an SM call while no enclaves are running.
The SM will first check that the new partitioning is valid, doesn't affect regions assigned to the SM.
It then waits for other cores to flush their L1 caches and to enter a blocked state in the SM, making sure no requests will be sent to the modified regions during reconfiguration (which could put the cache in an incoherent state).
The SM will then interact with the hardware to flush the LLC slices that are modified and set the LLC in its new configuration.

At the hardware level, we now design a practical reconfigurable LLC partitioning scheme.
We cannot use way partitioning \cite{kiriansky2018dawg} as the LLC has only 16 ways, and we need up to 64 slices, and Ironhide's \cite{ironhide} strategy assumes limited reconfiguration.
We propose a new, simple, configurable set partitioning scheme that only requires minimal changes in the hardware.

The configuration for each LLC slice (base and bound of the set-index range) is stored in a big register (1216 bits) called the set-index range table.
The SM can modify the entries of the set-index range table by sending MMIO requests to the cache, which is seen as a new device.
When the LLC receives a request, it will compute the new set index for the corresponding address.
We represent the transformation used in Figure \ref{fig:dynamicpartitioning}.
First, we extract a region ID from the address upper bits and obtain the corresponding index range base and size by looking up the set-index range table \circled{1}.
We then extract the address original set index, located in the middle of the address, and tweak it.
We compute the original set-index value modulo the size of the index range \circled{2} and add the base index \circled{3}.
This effectively restricts the new set index to the corresponding index range, partitioning the LLC \circled{4}.
Note that we must extend the length of the tags used by the LLC to include the entire original index in case an index range is of size 1 and all addresses from the same region map to the same set.
 
We observe on our cryptographic library that we reduce overheads from 120\% to 19.4\% by using a non-uniform partition for the LLC. More details can be found in Section \ref{sec:evalllcpartitioning}.

\begin{figure}[t]
\centering
\includegraphics[width=\columnwidth]{figures/partitioning.png}
\caption{LLC set indexing scheme. The function takes an address as an input and output a set index in the LLC slice allocated to its memory region.}
\label{fig:dynamicpartitioning}
\end{figure}

\noindent\textbf{LLC fine-grain flush}
To go hand-in-hand with our reconfigurable LLC partitioning, we also want to be able to selectively flush parts of the LLC.
We create a \emph{Zero Device} in the upper physical address space. 
It is implemented as a simple piece of logic connected to the cache hierarchy at the same level as the DRAM, but that always sends back zeros.
The trick is that the address space covered by the zero device has the same size as the DRAM and is located at addresses that will collide with DRAM addresses in the cache.
Once the SM has placed the machine in a correct waiting state (where all the cores have flushed their L1 and are waiting in the SM), it only needs to access an eviction set in the zero-device address space to flush the cache slice.
This means that flushing a small LLC slice is much faster than flushing a big one.
We look at the performance of our flushing mechanism in the evaluation section (see \ref{sec:evalllcpartitioning}).

\jules{Finish pass here}
\subsection{Pipelined Shared Accesses}

At this step, we observe that our test benchmark still present a 19.4\% overhead compared to our insecure benchmark.
We decide to measure the performance of different sections of the code to find and analyze the bottleneck.
We identify the function creating the slowdown to be located inside the SHA256 library.
This function performs burst accesses to shared memory to load the message to be signed.
This behavior is exacerbated by the compiler, trying to optimize the code by placing all the load instructions to shared memory one after the other.
Indeed, with our current mechanism, memory accesses to shared regions, first cannot use the L1 and have a 10 cycle latency, second cannot be pipelined as each instruction needs to wait to reach the head of the ROB to be sent.
This is especially costly during burst accesses. 
The LLC supporting 16 outstanding requests and having a 10 cycles latency, 16 memory accesses should be pipelined to take a total of $10 + 16 = 26$ cycles. Without pipelining, they will now take $10 * 16 = 160$ cycles.

As a result, we designed a mode where we temporarily change the speculation model to allow shared memory burst accesses (as pipelining is a form of speculation).
Nevertheless, we don't want to do so while enabling any type of speculation.
Indeed, we don't want our program to speculatively jump to a spectre-like gadget and transmit secrets to our attacker.
To avoid this, we tweak the speculation model to allow pipelined accesses to shared memory but disable the use and training of any predictors that could result in the speculative execution of rogue code.
Specifically, we disable the Branch Target Buffer (BTB), the Branch History Table (BHT) and the Return Address Stack (RAS).
Effectively, that guarantees that the hardware will always predict that the next program counter to be ``$pc+4$'' and avoid speculative execution of arbitrary code.
We call this \emph{straight-line} speculation.

Modifying the speculative model also modifies the hardware software contract.
The programmer now needs to ensure that when this mode is enabled, no speculative transmitter can be reached by the program when doing straight-line speculation.
Static analysis can be done manually or by reusing many existing tools: e.g. oo7\cite{wang2019oo7}, Kasper\cite{johannesmeyer2022kasper}, Kleespectre\cite{wang2020kleespectre}, Spectector\cite{guarnieri2020spectector}.
Note that the usual critique and limitation of these tools are that 1) the code bases are too big to analyze without producing a lot of false positives 2) they make simplifying assumptions on the speculation model.
Interestingly, both of those limitations are not concerns in the context of Citadel's burst mode because the code snippets we analyze are small (e.g. boiling down to a memcpy function for instance) and the speculation model is well-defined, simple and enforced (only speculates PC+4). Hence, these works are well-suited to our needs.
We consider adapting these tools to our more simple setup as future work.
We add two new configurations in our software-exposed configuration registers to enable and disable this ``burst'' mode.
The right way to use this mode is, like in our case, during burst accesses to shared memory (e.g. a memcpy function moving data from public memory to private enclave memory).
The programmer only needs to identify these series of burst access and guard them with instructions activating and deactivating burst mode.
When doing so on our cryptographic benchmark, we observe that overhead falls from $19.4\%$ to $2.7\%$.
More results are presented in the evaluation section.


\subsection{Relaxed Flushing Policies}
Because we build the entire software infrastructure required to run enclaves, we were able to observe that some security domains, especially the enclave and some parts of the mini-SM, do not have secrets to keep from each other.
Specifically, the mini-SM never manipulates sensitive data.
As a result, the mini-SM does not need to flush microarchitectural state when context switching to enclaves.
This reduces the cost of mini-SM call and opens the door to the use of frequent software or timer interrupts.
To disable flushing selectively, we modify the hardware to deactivate systematic flushing and exposed to the SM a register to trigger the flush.
We write the SM to trigger the flushing mechanism on required context switches, but not on most traps and SM-calls from the enclave.
We evaluate the performance implication of our mechanism in Section \ref{sec:evalrealxedflush}.

